% !TEX root = ../main.tex
\section{Discussion}
\label{sec:discussion}

\paragraph{Using the report}
Read one exam at a time, the report answers separate questions that a single score merges. Whether a
policy is safe because it yields or because it stays short is answered by exposure together with
hazard sensitivity: DiffusionDrive keeps its distance by planning short trajectories, which holds only
as long as the trajectories stay short, whereas AutoVLA and Alpamayo-1.5 cover a human-like distance
without the response to match. Whether a plan that moves is a plan that yields is answered by the
decomposition of \cref{sec:axes}: among the cells in which the pedestrian-blind plan does reach the
pedestrian, only \NumCfRealPct\% meet all three conditions, and displacement alone would have ranked the
least safe policy first.
The exam that no current benchmark reports is scaling: whether the response grows when the hazard
does. Five policies fail it, and the one that passes, DiffusionDriveV2, does so at a sensitivity below
0.15.

\paragraph{Limitations}
The exams are open-loop. The approach sequences are replayed from logs and re-anchored to the logged
pose. The night is a synthetic rendering that also halves the number of other pedestrians the detector
finds. The counterfactual pairs are constructed by editing, not captured, and so carry noise: 8.5\% of
removals leave a detectable residue, the inpainting artefact is 5--17\% of the effect, and the VLA
models can be edited only through their front camera. Alpamayo-1.5 covers 170 of the 236 frames. SimLingo plans 2.5\,s; for NAVSIM's 4\,s we extrapolate at constant velocity (holding position instead gives EPDMS 0.47, still the lowest).
Cross-policy correlations rest on six policies.

\paragraph{Future work}
The check-up is designed as the first half of a benchmark-to-deployment pairing, and as one hazard
type among several a deployment decision must cover: extending its exams to lead-vehicle braking and
turning or intersection conflicts is the natural next scenario. Within the pedestrian scenario, the
next step is to \emph{collect} paired data rather than construct it: the same scene with and without a
pedestrian, by day and by night. We then plan to validate the diagnoses against the same policies' behaviour on a real
vehicle. The sample-size laws above tell us how large that collection must be. A few dozen pairs
suffice for planned distance, over-reaction and lighting; hundreds of pairs with a genuine need to
react are needed for hazard sensitivity.

\section{Conclusion}
A benchmark score reports an outcome, not the behaviour behind it, so a rank cannot be trusted to
hold at a new site. We proposed a check-up that reads that behaviour from a few hundred logged frames
through verified, pixel-level counterfactual edits, reducing each scene to two axes of causal
attribution---the faithfulness displacement $F$ and the interference displacement $I$---on which five
exams are built. Three results carry the paper. First, on the cells in which a pedestrian does lie on
the planned path, only \NumCfRealPct\% of responses are genuine avoidance, and no policy either
diverts or slows; a photometric night-style perturbation moves every plan further than the pedestrian
does. Second, diagnosed on one driving side and tested on the other, the exam that transfers predicts
the new side's near-miss rate at $\rho={\NumSpBeh}$, where the standard scores reach only
$-0.54$--$+0.49$; what carries over is the ordering and the verdict. Third, each verdict can be priced
in advance---the lighting verdict within \NumNstarCFRhi{} frames, exposure within \NumNstarExpHi---so
that a deployment decision costs tens of target-domain frames rather than a fleet test.
